\documentclass[11pt,a4paper]{article}
\usepackage[utf8]{inputenc}
\usepackage[T1]{fontenc}
\usepackage[margin=0.75in]{geometry}
\usepackage{enumitem}
\usepackage{titlesec}
\usepackage{xcolor}
\usepackage{hyperref}

\definecolor{primarycolor}{RGB}{33,37,41}
\definecolor{secondarycolor}{RGB}{0,102,204}
\definecolor{accentcolor}{RGB}{108,117,125}
\hypersetup{colorlinks=true,linkcolor=secondarycolor,urlcolor=secondarycolor,citecolor=secondarycolor}
\titleformat{\section}{\Large\bfseries\color{primarycolor}}{}{0em}{}
\titlespacing*{\section}{0pt}{1.1em}{0.45em}
\setlist[itemize]{leftmargin=1.3em,itemsep=0.2em,topsep=0.25em}
\pagestyle{empty}
\newcommand{\cvrole}[3]{\noindent\textbf{#1} \hfill \textit{#2}\\[-0.15em]#3\par}

\begin{document}
\begin{center}
{\Huge\bfseries\color{primarycolor} Tobias Mwangi-Hartley}\\[0.35em]
{\large\color{accentcolor} Data Scientist | Applied Machine Learning}
\end{center}
\vspace{0.4em}

\section*{PROFESSIONAL SUMMARY}
Data scientist with experience developing forecasting and churn models and building experimentation tools. Hands-on with Python, BigQuery, Vertex AI Workbench, scheduled notebook workflows and weekly batch scoring. Interested in applying this experience while developing further expertise in the production operation of machine learning systems.

\section*{PROFESSIONAL EXPERIENCE}
\cvrole{Data Scientist | Thistlebox}{2022--Present}{}
\begin{itemize}
\item Developed a weekly six-week recipe demand forecast using LightGBM and calendar, price and menu-position features. WAPE improved from 23\% to 14\%, with an estimated food-waste reduction of 9\%; scheduled notebook runs write results to BigQuery for operations to use in Looker.
\item Built an XGBoost churn propensity model scored weekly in BigQuery; retention offers are sent to top-decile customers. A 10\% holdout validation showed churn down 2.1 percentage points in the treated group.
\item Built a Streamlit calculator for planning and analysing sequential A/B tests, used by around 25 people in product and growth. Wrote the app and its Dockerfile; the data platform team handles its deployment to Cloud Run.
\item Began writing pytest unit tests for forecasting feature code in 2025; coverage is approximately 60\%.
\item Delivered the internal training ``Stats for PMs'' three times.
\end{itemize}

\cvrole{Junior Data Analyst | Firth Energy Retail}{2020--2022}{}
\begin{itemize}
\item Segmented customers and analysed tariffs using R and SQL.
\item Automated monthly regulator reports with R Markdown, saving the team about two days each month.
\end{itemize}

\section*{SELECTED PROJECT}
\noindent\textbf{Image Classification Course Project | Personal | 2025}
\begin{itemize}
\item Fine-tuned a small CNN in PyTorch on a public food images dataset, achieving approximately 80\% accuracy. Learning project completed as part of an online deep learning course.
\end{itemize}

\section*{EDUCATION}
\noindent\textbf{MSc in Statistics}\\
Dissertation: Bayesian hierarchical models

\section*{TECHNICAL SKILLS}
\textbf{Programming and analysis:} Python, pandas, scikit-learn, R, SQL\\
\textbf{Machine learning and statistics:} LightGBM, XGBoost, SHAP, statsmodels, PyTorch, torchvision, feature engineering, forecasting, churn modelling, experimentation\\
\textbf{Data and development tools:} BigQuery, Vertex AI Workbench, dbt, Looker, Streamlit, Docker, pytest, R Markdown, Google Colab\\
\textbf{Workflows:} Scheduled notebook jobs, weekly batch scoring
\end{document}